\documentclass[manuscript,review,nonacm]{acmart}
%
% COMP5111 Assignment 2 -- Task 4 LLM-use report
%

\usepackage{listings}
\usepackage{xcolor}
\usepackage{booktabs}

\lstset{
  basicstyle=\ttfamily\small,
  breaklines=true,
  columns=fullflexible,
  frame=single,
  framesep=4pt,
  backgroundcolor=\color{gray!8},
  keywordstyle=\bfseries\color{blue!60!black},
  commentstyle=\itshape\color{green!50!black},
  stringstyle=\color{red!50!black},
  showstringspaces=false,
}

\setcopyright{none}
\settopmatter{printacmref=false, printccs=false, printfolios=true}
\renewcommand{\footnotetextcopyrightpermission}[1]{}
\pagestyle{plain}

\title{Task 4 \textendash{} Test-Driven LLM Code Generation\\
       \large COMP5111 Assignment 2}
\author{Zimo Ji}
\affiliation{\institution{HKUST}\city{Hong Kong}\country{Hong Kong SAR}}

\begin{document}
\maketitle
\thispagestyle{plain}

\noindent\textit{The companion \texttt{Task4\_Iteration\_Trace.pdf}
has the verbatim prompt + response + extracted code for every round.
This report is the analysis only.}

\section{Setup}

\subsection{Selected methods}
I picked ten public-static methods from \texttt{Subject.java} and put
them in a reduced CUT (\texttt{reports/Task4\_Data/NewSubject.java}),
plus one private helper (\texttt{isOneOf}) needed by
\texttt{parseToken}. All public signatures are byte-identical to the
originals:

\begin{small}
\begin{tabular}{rll}
\toprule
\# & Signature & From \texttt{Subject.java} \\
\midrule
1  & \texttt{boolean startsWithIgnoreCase(String, String)} & \texttt{StringAlgorithms} \\
2  & \texttt{String  parseToken(String, char[])}           & \texttt{StringAlgorithms} \\
3  & \texttt{int     extractIntInStr(String)}              & \texttt{StringAlgorithms} \\
4  & \texttt{int[]   getVersionNo(String)}                 & \texttt{StringAlgorithms} \\
5  & \texttt{String  padLeft(String, short, char)}         & \texttt{StringAlgorithms} \\
6  & \texttt{boolean judgeLeapYear(int)}                   & \texttt{DateTimeAlgorithms} \\
7  & \texttt{int     calcDaysInMonth(int, int)}            & \texttt{DateTimeAlgorithms} \\
8  & \texttt{int     getQuarter(int)}                      & \texttt{DateTimeAlgorithms} \\
9  & \texttt{int     monAbbr2month(String)}                & \texttt{DateTimeAlgorithms} \\
10 & \texttt{String  month2MonAbbr(int)}                   & \texttt{DateTimeAlgorithms} \\
\bottomrule
\end{tabular}
\end{small}

I used the Task 2 fault-fixed implementations as ground truth so that
EvoSuite's tests pin down the \emph{intended} contract, not the seeded
faults.

\subsection{EvoSuite generation}

\begin{lstlisting}
java -jar lib/evosuite-1.2.0.jar \
    -class comp5111.assignment.cut.NewSubject \
    -projectCP reports/Task4_Data/build/classes \
    -Dsearch_budget=600 -Dstopping_condition=MaxTime \
    -Dcriterion=LINE:BRANCH -Dassertion_strategy=ALL \
    -Drandom_seed=2026
\end{lstlisting}

100\,\% LINE coverage (115/115) and 100\,\% BRANCH coverage (131/131)
in 1\,s of search; 62 JUnit tests in the 414-line
\texttt{NewSubject\_ESTest.java}.

\subsection{LLM endpoint}
OpenAI-compatible \texttt{/chat/completions},
\texttt{temperature=0.1}. Primary model is \texttt{gpt-5.1}; I also
ran \texttt{gpt-5.4} (stronger, footnoted in \S3) and
\texttt{gpt-4o-mini} (weaker, ablation in \S6). Raw HTTP from the
Python standard library, no SDK.

\section{Prompt design}

The same template every round. Four fixed parts; rounds $\ge$ 2 add
two iteration-specific parts:

\begin{lstlisting}
<system>
You are a senior Java engineer. Respond with exactly one fenced ```java code
block as instructed.
</system>

<user>
You are a Java engineer. Implement the class
`comp5111.assignment.cut.NewSubject` so that every JUnit test in the
provided test suite passes.

Rules:
 - Output exactly ONE fenced ```java code block containing the FULL contents
   of NewSubject.java (package, imports if any, the entire class).
 - Do NOT modify the test suite. Do NOT add main(). Keep the package and
   class name unchanged. Method signatures shown below are mandatory.
 - Output nothing outside the single code block.

=== Method signatures + Javadoc (mandatory shape) ===
<NewSubject.java with method bodies stripped to ";">

=== EvoSuite-generated test suite (passes against the ground truth) ===
<entire NewSubject_ESTest.java verbatim>

=== Your previous implementation (round N-1) ===           # rounds >= 2
<last LLM output>

=== Test failures observed for that implementation ===     # rounds >= 2
<list of "FAILURE>>>...<<<FAILURE" lines>
</user>
\end{lstlisting}

I kept the Javadoc on purpose: it carries intent (e.g.\
``\texttt{extractIntInStr} returns the LAST contiguous digit run'')
that the test suite samples but does not state. Round-1 prompt is
15\,138 chars; round 2 grows to about 24\,000 once I append the
previous attempt and failing tests.

\section{First-round result}

\begin{tabular}{lc}
\toprule
Metric          & Round 1 \\
\midrule
Compilation OK  & yes \\
Tests run       & 62 \\
Tests failed    & \textbf{1} \\
Pass rate       & \textbf{98.4\,\% (61/62)} \\
\bottomrule
\end{tabular}

\medskip

The single failure was

\begin{lstlisting}
test51 :: arrays first differed at element [0];
         expected:<6> but was:<316>
\end{lstlisting}

\texttt{test51} calls
\texttt{getVersionNo("\$S3P\&,*V+D\&1\{\%U6")} and expects
\texttt{[6,0,0,0]}. Input has no \texttt{.}, so
\texttt{String.split("\textbackslash{}\textbackslash{}.", -1)} returns
one segment. \texttt{gpt-5.1} walked the segment forward and
concatenated \emph{every} digit (\texttt{3}, \texttt{1}, \texttt{6}
$\to$ \texttt{316}) instead of returning only the last contiguous
digit run as the Javadoc says (\texttt{"\dots\&1\{\%U6"} $\to$
\texttt{6}). The other 61 tests don't exercise this case (they have
a leading digit, so forward and backward scans agree, or no digits at
all).

\medskip
\noindent\textbf{Footnote (stronger backend).} With \texttt{gpt-5.4}
and the same prompt, all 62 tests pass on round 1, no iteration
needed. Artefacts under \texttt{reports/Task4\_Data/trial\_gpt54/}. I
do not use it as the primary experiment because it leaves the
iteration loop with nothing to do.

\section{Final pass rate after refinement}

I appended round 1's failure (\texttt{expected:<6> but was:<316>}) to
the prompt. Round 2 fixed it: \texttt{gpt-5.1} rewrote the digit loop
in \texttt{getVersionNo} to scan \textbf{backwards} from the end of
the segment and record the last contiguous digit run. The diff inside
\texttt{getVersionNo}:

\begin{lstlisting}[language=Java]
// round 1 (buggy): forward accumulator
for (int j = 0; j < p.length(); j++) {
    char c = p.charAt(j);
    if (Character.isDigit(c)) { hasDigit = true;
        num = num * 10 + (c - '0');
    }
}

// round 2 (fixed): backward search for the last digit run
int end = -1, start = -1;
for (int j = p.length() - 1; j >= 0; j--) {
    char c = p.charAt(j);
    if (Character.isDigit(c)) {
        if (end == -1) end = j;
        start = j;
    } else if (end != -1) {
        break;
    }
}
int num = 0;
for (int j = start; j <= end; j++) num = num * 10 + (p.charAt(j) - '0');
\end{lstlisting}

Round 2 passes all 62 tests; my loop early-exits.

\medskip
\noindent\textbf{Final: 100\,\% (62/62) after 2 rounds.}

\medskip

Per-round artefacts:

\begin{small}
\begin{tabular}{ll}
\toprule
Path & Contents \\
\midrule
\texttt{reports/Task4\_Data/prompts/round\_\{1,2\}.txt}              & exact prompts                    \\
\texttt{reports/Task4\_Data/llm\_outputs/round\_\{1,2\}.txt}         & raw LLM responses                \\
\texttt{reports/Task4\_Data/generated/round\_\{1,2\}/NewSubject.java}& extracted Java each round        \\
\texttt{reports/Task4\_Data/test\_results/round\_\{1,2\}.txt}        & compile + JUnit summary (json)   \\
\texttt{reports/Task4\_Data/final/NewSubject.java}                   & round-2 implementation (= final) \\
\texttt{reports/Task4\_Data/summary.json}                            & per-round metrics                \\
\texttt{reports/LLM/Task4\_Iteration\_Trace.\{md,pdf\}}              & full trace                       \\
\bottomrule
\end{tabular}
\end{small}

\section{Similarities and differences vs Subject.java}

I diffed the \textbf{final (round 2)} \texttt{NewSubject.java} from
\texttt{gpt-5.1} against the matching ten methods of
\texttt{Subject.java}.

\subsection{Identical (modulo whitespace / comments)}
\texttt{judgeLeapYear}, \texttt{calcDaysInMonth},
\texttt{getQuarter}, \texttt{month2MonAbbr}, \texttt{padLeft}. Same
control flow, same constants.

\subsection{Same intent, different idiom}
\begin{itemize}\itemsep1pt
\item \texttt{startsWithIgnoreCase}. Original copies
      \texttt{prefix.length()} chars and compares;
      \texttt{gpt-5.1} walks the prefix with
      \texttt{Character.toLowerCase(c1) != Character.toLowerCase(c2)}.
      The LLM's loop avoids the substring allocation, behaviour is the
      same.
\item \texttt{getVersionNo}. Original walks with
      \texttt{indexOf('.')} and a manual cursor in a
      \texttt{do\,\dots\,while}; \texttt{gpt-5.1} (round 2) uses
      \texttt{versionString.split("\textbackslash{}\textbackslash{}.", -1)},
      validates \texttt{parts.length} and empty parts, then runs the
      backward digit-run scan above. Same \texttt{int[4]} output.
\item \texttt{parseToken}. Original returns
      \texttt{str.substring(0, pos)} after a \texttt{while} loop;
      \texttt{gpt-5.1} adds null/empty checks and returns
      \texttt{str.substring(0, i)} directly inside the
      matched-terminator branch. The buggy \texttt{pos++} of
      \texttt{Subject.java} line~205 is absent (EvoSuite tested the
      patched ground truth).
\item \texttt{extractIntInStr}. Original is a single forward pass
      that resets the accumulator on every non-digit.
      \texttt{gpt-5.1} does a backward scan for the last digit run,
      then a small forward accumulator inside it. Same contract, more
      verbose.
\end{itemize}

\subsection{Different control structure}
\texttt{monAbbr2month}. Original uses a perfect hash
\texttt{(ch0<<16) | (ch1<<8) | ch2} and dispatches on twelve
hard-coded integers -- one of which (\texttt{5465466} for ``Sep'')
is the seeded fault from Task~2. \texttt{gpt-5.1} drops the hash and
uses a Java~7+ \texttt{switch (abrr)}:

\begin{lstlisting}[language=Java]
switch (abrr) { case "Jan": return 1; ... case "Sep": return 9; ... }
\end{lstlisting}

Same return values, but the perfect-hash trick is gone. The LLM
couldn't have known the original used it -- EvoSuite tests only
observe return values.

\subsection{LLM-introduced defensive checks}
\texttt{gpt-5.1} consistently adds null/empty pre-checks the original
sometimes omits (e.g.\ \texttt{parseToken} against null \texttt{str}
or zero-length \texttt{terminators}). Doesn't break tests here, but
would change the contract if the spec required NPE on null. The same
instinct goes wrong on a weaker model (\S6).

\section{Ablation: weaker backend model}

To get iteration data and check how brittle round-1 convergence is, I
re-ran the same prompt template, EvoSuite test suite, and
\texttt{NewSubject.java} with \texttt{gpt-4o-mini}. Everything else
(\texttt{temperature=0.1}, \texttt{max-rounds=5}, OpenAI-compatible
endpoint) is identical.

\begin{tabular}{ccccc}
\toprule
Round & compile\_ok & run & fail & pass rate \\
\midrule
1 & yes & 62 & 2 & 96.8\,\% \\
2 & yes & 62 & 2 & 96.8\,\% \\
3 & yes & 62 & 2 & 96.8\,\% \\
4 & yes & 62 & 2 & 96.8\,\% \\
5 & yes & 62 & 2 & 96.8\,\% \\
\bottomrule
\end{tabular}

\medskip

The weaker model converges fast to a near-correct implementation but
\textbf{doesn't fix the last two failures} even after four rounds of
feedback:

\begin{itemize}\itemsep1pt
\item \texttt{test43}: \texttt{calcDaysInMonth(-3324, 2)} expected
      \texttt{29}, got \texttt{28}.
\item \texttt{test51}:
      \texttt{getVersionNo("\$S3P\&,*V+D\&1\{\%U6")} expected
      \texttt{[6,0,0,0]}, got \texttt{null}
      (\texttt{assertNotNull} failure).
\end{itemize}

Both come from over-defensive code that the model wrote on round 1
and kept every subsequent round:

\begin{itemize}\itemsep1pt
\item In \texttt{judgeLeapYear}, \texttt{gpt-4o-mini} added an
      unjustified guard
      \texttt{if (year < 0) return false;}. The Javadoc says nothing
      about the sign; the Gregorian rule works fine for negative
      years. With this guard, every negative leap year is
      mis-classified, so February of $-3324$ returns 28.
\item In \texttt{getVersionNo}, \texttt{gpt-4o-mini} called
      \texttt{Integer.parseInt(parts[i])} directly instead of the
      spec's ``last digit run'' extraction (the original calls
      \texttt{extractIntInStr}). \texttt{"\$S3P\&,*V+D\&1\{\%U6"} has
      no \texttt{.} and no leading digit, so \texttt{parseInt} throws
      \texttt{NumberFormatException}, and the LLM's
      \texttt{catch (NumberFormatException e) \{ return null; \}}
      returns \texttt{null} instead of \texttt{[6,0,0,0]}.
\end{itemize}

The model received the failure messages verbatim every round. Round 4
sent it:

\begin{lstlisting}
2 of 62 tests failed.
 - test43(...) :: expected:<29> but was:<28>
 - test51(...) :: null
\end{lstlisting}

\dots and \texttt{gpt-4o-mini} still produced ``fixes'' with the same
two bugs. The trace shows it rewriting the surrounding code each
round (different idioms, slightly different control flow) but leaving
the two over-defensive guards untouched. Two takeaways:

\begin{enumerate}\itemsep1pt
\item \textbf{Localising a bug from a high-level test message is hard
      for weak models.} ``expected 29, was 28'' doesn't say which
      line. Working out that the bug is in \texttt{judgeLeapYear}
      (not \texttt{calcDaysInMonth}) takes counterfactual reasoning
      that \texttt{gpt-4o-mini} doesn't do reliably.
\item \textbf{The model has a strong prior toward defensive code}
      that the spec doesn't require, and feedback that the defence is
      wrong is not enough to override that prior.
\end{enumerate}

\texttt{gpt-5.4} (\S3 footnote) writes none of these guards and
passes round 1 outright; primary \texttt{gpt-5.1} needs one round of
feedback. Same prompt, same Javadoc, same tests; the difference is
purely model capability. Per-round traces are in
\texttt{reports/LLM/Task4\_Iteration\_Trace\_weak\_gpt4o\_mini.\{md,pdf\}}
and \texttt{reports/Task4\_Data/weak\_gpt4o\_mini/}.

\section{Insights}

\textbf{Strengths.}
\begin{itemize}\itemsep1pt
\item \emph{Synthesis from spec + tests is fast and accurate.} On a
      $\approx$250 LOC CUT with Javadoc + 100\,\% coverage tests,
      \texttt{gpt-5.1} got 9 of 10 methods right on round 1 and the
      last one after one round of feedback (\S4). \texttt{gpt-5.4}
      passes round 1. Almost no prompt engineering beyond format
      constraints.
\item \emph{Iteration can localise and fix bugs from a high-level
      test message.} In round 2, \texttt{gpt-5.1} took
      \texttt{expected:<6> but was:<316>} and correctly diagnosed the
      bug in \texttt{getVersionNo}'s digit extraction (not the test,
      not \texttt{extractIntInStr}, not \texttt{String.split}).
      Non-trivial inference.
\item \emph{Idiomatic API choice.} Modern Java
      (\texttt{switch (String)}, \texttt{String.split},
      \texttt{Character.toLowerCase} loops) where legacy code would
      not.
\end{itemize}

\textbf{Limitations.}
\begin{itemize}\itemsep1pt
\item \emph{Hidden invariants are erased.} The perfect-hash dispatch
      in \texttt{monAbbr2month} is gone. A caller that depended on
      the exact integer constants would break silently. Tests don't
      observe internal control flow, so test-driven generation can't
      recover it.
\item \emph{No ``minimum-change'' notion.} Asked to reimplement, the
      LLM rewrites. For single-line patching (Task 2 elsewhere),
      this is the wrong tool.
\item \emph{Spec-test conflict goes to the test.} If the Javadoc and
      the EvoSuite tests disagree, the LLM follows the tests and
      ignores the doc.
\item \emph{Iteration doesn't always fix things (\S6).} On
      \texttt{gpt-4o-mini} the same two bugs survived four rounds of
      failure feedback.
\end{itemize}

\textbf{Limitations of EvoSuite-generated test cases.}
\begin{itemize}\itemsep1pt
\item \emph{Coverage saturates fast on simple methods.} 100\,\%
      LINE+BRANCH in one second, but the tests are mostly
      value-instance assertions
      (\texttt{assertEquals("Sep", month2MonAbbr(9))}) that don't
      express the contract.
\item \emph{No equivalence-class abstraction.} 12 separate tests for
      \texttt{month2MonAbbr} (one per valid month) plus several
      invalid; a hand-written test would be parametric. The bloat
      dominates the prompt budget (the test suite is over 95\,\% of
      round-1 prompt size).
\item \emph{Tests reflect whatever the implementation does.} If the
      ground truth had been buggy, EvoSuite would have locked the bug
      into the tests, and the LLM would have reproduced it.
\end{itemize}

\textbf{Suggestions for improving success rate.}
\begin{enumerate}\itemsep1pt
\item Always include the Javadoc, not just the tests. This was the
      biggest lever in my run; without it the perfect-hash dispatch
      in \texttt{monAbbr2month} could easily be reproduced as a
      buggy table.
\item In iteration rounds, send only the failing-test bodies + a
      diff of the previous attempt, not the whole suite. My 11\,KB
      suite is already heavy; on bigger CUTs it would exceed the
      model's effective context.
\item For single-line patches, use a different prompt shape -- full
      method, suspected line(s) highlighted, failing test, and an
      explicit ``do not rewrite any other line''.
\item Augment EvoSuite with property-based tests (jqwik etc.).
      Properties like \texttt{monAbbr2month(month2MonAbbr(m)) == m}
      for $1\le m \le 12$ capture invariants that examples can't.
\item Add a ``minimal-diff'' objective in the system message when
      refactoring or patching.
\end{enumerate}

\noindent\textit{(Body word count, excluding code: about 1100.)}

\end{document}
